\section{Ôn tập chương 3-Các số đặc trưng đo xu thế trung tâm cho mẫu số liệu ghép nhóm}
\Opensolutionfile{ans}[ans/ans-OTC3-TN]
\caulc
\begin{ex}%[1D5N1-2]%Câu 1
Điều tra về chiều cao của học sinh khối lớp 11, ta được mẫu số liệu sau:
\begin{center}
\begin{tabular}{|c|c|}
\hline Chiều cao $(\mathrm{cm})$ & Số học sinh \\
\hline$[150; 152)$ & $5$ \\
\hline$[152; 154)$ & $18$ \\
\hline$[154; 156)$ & $40$ \\
\hline$[156; 158)$ & $26$ \\
\hline$[158; 160)$ & $8$ \\
\hline$[160; 162)$ & $3$ \\
\hline Tổng & $N=100$ \\
\hline
\end{tabular}
\end{center}
Mẫu số liệu ghép nhóm đã cho có tất cả bao nhiêu nhóm?
	\choice
	{$5$}
	{\True $6$}
	{$7$}
	{$12$}
	\loigiai{
Mẫu số liệu ghép nhóm đã cho có tất cả $6$ nhóm.
}
\end{ex}

\begin{ex}%[1D5N1-2]%Câu 2
Cho mẫu số liệu về chiều cao $(\mathrm{cm})$ của các học $\sinh$ nữ trong khối 11 của một trường như sau:
\begin{center}
\begin{tabular}{|l|c|c|c|c|c|c|}
\hline Chiều cao & {$[145; 150)$} & {$[150; 155)$} & {$[155; 160)$} & {$[160; 165)$} & {$[165; 170)$} & {$[170; 175)$} \\
\hline Số học sinh & $20$ & $45$ & $34$ & $27$ & $15$ & $4$ \\
\hline
\end{tabular}
\end{center}
Mẫu số liệu trên có bao nhiêu số liệu, bao nhiêu nhóm?
	\choice
	{$145$ số liệu; $6$ nhóm}
	{$30$ số liệu; $5$ nhóm}
	{\True $6$ số liệu; $145$ nhóm}
	{$5$ số liệu; $30$ nhóm}
	\loigiai{
Mẫu số liệu $(T)$ có
$$20+45+34+27+15+4=145\text{ (số liệu).}$$
$6$ nhóm là $[145; 150)$; $[150; 155)$; $[155; 160)$; $[160; 165)$; $[165; 170)$; $[170; 175)$.
}
\end{ex}

\begin{ex}%[1D5N1-2]%Câu 3
Tìm hiểu thời gia xem tivi trong tuần trước (đơn vị: giờ) của một số học sinh thu được kết quả sau:
\begin{center}
\begin{tabular}{|l|c|c|c|c|c|}
\hline Thời gian (giờ) & {$[0; 5)$} & {$[5; 10)$} & {$[10; 15)$} & {$[15; 20)$} & {$[20; 25)$} \\
\hline Số học sinh & $8$ & $16$ & $4$ & $2$ & $2$ \\
\hline
\end{tabular}
\end{center}
Giá trị đại diện của nhóm $[20; 25)$ là
	\choice
	{\True $22{,}5$}
	{$23$}
	{$20$}
	{$5$}
	\loigiai{
Giá trị đại diện của nhóm $[20; 25)$ là $\dfrac{20+25}{2}=22{,}5$.
}
\end{ex}

\begin{ex}%[1D5N1-4]%Câu 4
Khảo sát thời gian tập thể dục của một số học sinh khối 11 thu được mẫu số liệu ghép nhóm sau:
\begin{center}
\begin{tabular}{|l|c|c|c|c|c|}
\hline Thời gian (phút) & {$[0; 20)$} & {$[20; 40)$} & {$[40; 60)$} & {$[60; 80)$} & {$[80; 100)$} \\
\hline Số học sinh & $5$ & $9$ & $12$ & $10$ & $6$\\
\hline
\end{tabular}
\end{center}
Nhóm chứa mốt của mẫu số liệu trên là
	\choice
	{\True $[40; 60)$}
	{$[20; 40)$}
	{$[60; 80)$}
	{$[80; 100)$}
	\loigiai{
Mốt $M_0$ chứa trong nhóm $[40; 60)$.
}
\end{ex}

\begin{ex}%[1D5H2-2]%Câu 5
Khảo sát thời gian tập thể dục của một số học sinh khối 11 thu được mẫu số liệu ghép nhóm sau:
\begin{center}
\begin{tabular}{|l|c|c|c|c|c|}
\hline Thời gian (phút) & {$[0; 20)$} & {$[20; 40)$} & {$[40; 60)$} & {$[60; 80)$} & {$[80; 100)$} \\
\hline Số học sinh & $5$ & $9$ & $12$ & $10$ & $6$ \\
\hline
\end{tabular}
\end{center}
Nhóm chứa trung vị của mẫu số liệu trên là
	\choice
	{\True $[40; 60)$}
	{$[20; 40)$}
	{$[60; 80)$}
	{$[80; 100)$}
	\loigiai{
Ta có $n=42$.\\
Nên trung vị của mẫu số liệu trên là $Q_2=\dfrac{x_{21}+x_{22}}{2}$.\\
Mà $x_{21}, x_{22} \in[40; 60)$ nên nhóm chứa trung vị của mẫu số liệu trên là nhóm $[40; 60)$.
}
\end{ex}

\begin{ex}%[1D5H2-3]%Câu 6
Khảo sát thời gian tập thể dục của một số học sinh khối 11 thu được mẫu số liệu ghép nhóm sau:
\begin{center}
\begin{tabular}{|l|c|c|c|c|c|}
\hline Thời gian (phút) & {$[0; 20)$} & {$[20; 40)$} & {$[40; 60)$} & {$[60; 80)$} & {$[80; 100)$} \\
\hline Số học sinh & $5$ & $9$ & $12$ & $10$ & $6$ \\
\hline
\end{tabular}
\end{center}
Nhóm chứa tứ phân vị thứ ba của mẫu số liệu trên là
	\choice
	{$[40; 60)$}
	{$[20; 40)$}
	{\True $[60; 80)$}
	{$[80; 100)$}
	\loigiai{
Ta có $n=42$.\\
Nên tứ phân vị thứ nhất của mẫu số liệu trên là $Q_3=x_{33}$.\\
Mà $x_{33} \in[60; 80)$ nên nhóm chứa tứ phân vị thứ nhất của mẫu số liệu trên là nhóm $[60; 80)$.
}
\end{ex}

\begin{ex}%[1D5H1-3]%Câu 7
Doanh thu bán hàng trong $20$ ngày được lựa chọn ngẫu nhiên của một của hàng được ghi lại ở bảng sau (đơn vị: triệu đồng).
\begin{center}
\begin{tabular}{|l|c|c|c|c|c|}
\hline Doanh thu & {$[5; 7)$} & {$[7; 9)$} & {$[9; 11)$} & {$[11; 13)$} & {$[13; 15)$} \\
\hline Số ngày & $2$ & $7$ & $7$ & $3$ & $1$ \\
\hline
\end{tabular}
\end{center}
Số trung bình của mẫu số liệu trên thuộc khoảng nào trong các khoảng dưới đây?
	\choice
	{$[7; 9)$}
	{\True $[9; 11)$}
	{$[11; 13)$}
	{$[13; 15)$}
	\loigiai{
Bảng tần số ghép nhóm theo giá trị đại diện là
\begin{center}
\begin{tabular}{|l|c|c|c|c|c|}
\hline Doanh thu & {$[5; 7)$} & {$[7; 9)$} & {$[9; 11)$} & {$[11; 13)$} & {$[13; 15)$} \\
\hline Giá trị đại diện & $6$ & $8$ & $10$ & $12$ & $14$ \\
\hline Số ngày & $2$ & $7$ & $7$ & $3$ & $1$ \\
\hline
\end{tabular}
\end{center}
Số trung bình là $\overline{x}=\dfrac{2\cdot6+7\cdot8+7\cdot10+3\cdot12+1\cdot14}{20}=9{,}4$.
}
\end{ex}

\begin{ex}%[1D5V2-2]%Câu 8
Cho mẫu số liệu ghép nhóm về thống kê điểm số (thang điểm $20$) của 100 học sinh tham dự kỳ thi học sinh giỏi toán, ta có bảng số liệu sau:
\begin{center}
\begin{tabular}{|l|c|c|c|c|c|c|}
\hline Điểm & {$[8; 10)$} & {$[10; 12)$} & {$[12; 14)$} & {$[14; 16)$} & {$[16; 18)$} & {$[18; 20)$} \\
\hline Số học sinh & $6$ & $21$ & $30$ & $25$ & $14$ & $4$\\
\hline
\end{tabular}
\end{center}
Tìm trung vị của mẫu số liệu ghép nhóm trên. (Kết quả làm tròn đến hàng phần trăm).
	\choice
	{$12{,}18$}
	{$12{,}81$}
	{$13{,}35$}
	{\True $13{,}53$}
	\loigiai{
Cỡ mẫu là $n=6+21+30+25+14+4=100$.\\
Gọi $x_1; x_2; \ldots; x_{100}$ là số điểm của 100 học sinh và được sắp xếp theo thứ tự tăng dần. Khi đó, trung vị là $\dfrac{x_{50}+x_{51}}{2}$, do hai giá trị $x_{50}, x_{51}$ thuộc nhóm $[12; 14)$ nên nhóm này chứa trung vị.\\
Ta có $p=3$; $a_3=12$; $m_3=30$; $m_1+m_2=25$; $a_4-a_3=2$ khi đó trung vị của mẫu số liệu là
$$M_e=12+\dfrac{\dfrac{100}{2}-27}{30}\cdot 2=13{,}53.$$
}
\end{ex}

\begin{ex}%[1D5N1-4]%Câu 9
Cho mẫu số liệu về thời gian (phút) đi từ nhà đến trường của các học sinh trong một lớp 11 của một trường như sau:
\begin{center}
\begin{tabular}{|l|c|c|c|c|c|c|}
\hline Thời gian & {$[0; 5)$} & {$[5; 10)$} & {$[10; 15)$} & {$[15; 20)$} & {$[20; 25)$} & {$[25; 30)$} \\
\hline Số học sinh & $7$ & $12$ & $7$ & $5$ & $3$ & $2$ \\
\hline
\end{tabular}
\end{center}
Có bao nhiêu học sinh có thời gian đi từ nhà đến trường là $15$ phút đến $20$ phút?
	\choice
	{$20$}
	{$15$}
	{\True $5$}
	{$7$}
	\loigiai{
Theo bảng số liệu, có $5$ học sinh có thời gian đi từ nhà đến trường là $15$ phút đến $20$ phút.
}
\end{ex}

\begin{ex}%[1D5H1-3]%Câu 10
Cho mẫu số liệu ghép nhóm về thống kê nhiệt độ tại một địa điểm trong $30$ ngày, ta có bảng số liệu sau:
\begin{center}
\begin{tabular}{|l|c|c|c|c|}
\hline Nhiệt độ $\left({ }^{\circ} \mathrm{C}\right)$ & {$[18; 21)$} & {$[21; 24)$} & {$[24; 27)$} & {$[27; 30)$}\\
\hline Số ngày & $6$ & $12$ & $9$ & $3$ \\
\hline
\end{tabular}
\end{center}
Nhiệt độ trung bình trong $30$ ngày trên là
	\choice
	{$24^{\circ} \mathrm{C}$}
	{$25{,}4^{\circ} \mathrm{C}$}
	{$24{,}3^{\circ} \mathrm{C}$}
	{\True $23{,}4^{\circ} \mathrm{C}$}
	\loigiai{
Giá trị đại diện của mỗi nhóm số liệu là trung bình cộng của hai đầu mút.\\
Ta có bảng tần số ghép nhóm theo giá trị đại diện của mỗi nhóm:
\begin{center}
\begin{tabular}{|l|c|c|c|c|}
\hline Nhóm & {$[18; 21)$} & {$[21; 24)$} & {$[24; 27)$} & {$[27; 30)$}\\
\hline Giá trị đại diện & $19{,}5$ & $22{,}5$ & $25{,}5$ & $28{,}5$ \\
\hline Tần số & $6$ & $12$ & $9$ & $3$ \\
\hline
\end{tabular}
\end{center}
Nhiệt độ trung bình trong 30 ngày trên là
$$\overline{x}=\dfrac{6 \cdot 19{,}5+12 \cdot 22{,}5+9\cdot25{,}5+3 \cdot 28{,}5}{30}=23{,}4^{\circ} \mathrm{C}.$$
}
\end{ex}

\begin{ex}%[1D5N1-4]%Câu 11
Cho mẫu số liệu về chiều cao $(\mathrm{cm})$ của các học sinh nữ trong khối 11 của một trường như sau:
\begin{center}
\begin{tabular}{|l|c|c|c|c|c|c|}
\hline Chiều cao & {$[145; 150)$} & {$[150; 155)$} & {$[155; 160)$} & {$[160; 165)$} & {$[165; 170)$} & {$[170; 175)$} \\
\hline Số học sinh & $20$ & $45$ & $34$ & $27$ & $15$ & $4$ \\
\hline
\end{tabular}
\end{center}
Số học sinh nữ cao từ $150 \mathrm{~cm}$ đến $155 \mathrm{~cm}$ là
	\choice
	{$20$}
	{$65$}
	{$34$}
	{\True $45$}
	\loigiai{
Theo bảng số liệu trên, số học sinh nữ cao từ $150 \mathrm{~cm}$ đến $155 \mathrm{~cm}$ là $45$ học sinh.
}
\end{ex}

\begin{ex}%[1D5V2-2]%Câu 12
Cho mẫu số liệu ghép nhóm về chiều cao của $25$ cây dừa giống như sau:
\begin{center}
\begin{tabular}{|l|c|c|c|c|c|}
\hline Chiều cao $(\mathrm{cm})$ & {$[0; 10)$} & {$[10; 20)$} & {$[20; 30)$} & {$[30; 40)$} & {$[40; 50)$} \\
\hline Số cây & $4$ & $6$ & $7$ & $5$ & $3$ \\
\hline
\end{tabular}
\end{center}
Trung vị của mẫu số liệu ghép nhóm này là
	\choice
	{$M_e=\dfrac{175}{7}$}
	{$M_e=\dfrac{165}{5}$}
	{$M_e=\dfrac{165}{7}$}
	{\True $M_e=\dfrac{165}{3}$}
	\loigiai{
Cỡ mẫu $n=4+6+7+5+3=25$.\\
$x_1, x_2, \ldots, x_{25}$ là chiều cao của $25$ cây dừa giống được sắp xếp theo thứ tự không giảm. Khi đó, trung vị là $x_{13}$. Do $x_{13}$ thuộc nhóm $[20; 30)$ nên nhóm này chứa trung vị.\\
 Do đó $p=3$, $a_3=20$, $m_3=7$, $m_1+m_2=10$, $a_4-a_3=10$.\\
Vậy $M_e=20+\dfrac{\dfrac{25}{2}-10}{7} \cdot 10=\dfrac{165}{7}$.
}
\end{ex}
\Closesolutionfile{ans}
\indapan{6}{ans/ans-OTC3-TN}

\cauds
\Opensolutionfile{ans}[ans/ans-OTC3-DS]
\setcounter{ex}{0}
\begin{ex}%[1D5H1-4]%Câu 1
Cho mẫu số liệu điểm môn Toán của một nhóm học sinh như sau:
\begin{center}
\begin{tabular}{|l|c|c|c|c|}
\hline Điểm & $[6; 7)$ & $[7; 8)$ & $[8; 9)$ & $[9; 10]$ \\
\hline Số học sinh & $8$ & $7$ & $10$ & $5$ \\
\hline
\end{tabular}
\end{center}
Xét tính đúng sai của các mệnh đề sau:
	\choiceTF
	{\True Mẫu số liệu đã cho là mẫu số liệu ghép nhóm}
	{\True Cỡ mẫu của mẫu số liệu là $30$}
	{\True Điểm trung bình của các học sinh là $7{,}9$}
	{Mốt của mẫu số liệu là $10$}
	\loigiai{
Ta có bảng sau
\begin{center}
\begin{tabular}{|l|c|c|c|c|}
\hline Điểm & {$[6; 7)$} & {$[7; 8)$} & {$[8; 9)$} & {$[9; 10]$} \\
\hline Giá trị đại diện & $6{,}5$ & $7{,}5$ & $8{,}5$ & $9{,}5$ \\
\hline Số học sinh & $8$ & $7$ & $10$ & $5$ \\
\hline
\end{tabular}
\end{center}
	\begin{itemchoice}
	\itemch Đúng. Vì theo định nghĩa mẫu số liệu ghép nhóm.
	\itemch Đúng. Vì cỡ mẫu là $8+7+10+5=30$.
	\itemch Đúng. Vì điểm trung bình của các học sinh là $$\overline{x}=\dfrac{6{,}58+7{,}5 \cdot 7+8{,}5 \cdot 10+9{,}5 \cdot 5}{30}=7{,}9.$$
	\itemch Sai. Vì nhóm chứa Mốt là $[8; 9)$.\\
 Mốt của mẫu số liệu là $M_e=8+\dfrac{10-7}{10-7+10-5}(9-8)=8{,}375$.
	\end{itemchoice}
}
\end{ex}

\begin{ex}%[1D5H2-3]%Câu 2
Cho mẫu số liệu ghép nhóm về lương của nhân viên trong một công ty như sau:
\begin{center}
\begin{tabular}{|l|c|c|c|c|c|}
\hline Lương (triệu đồng) & {$[9; 12)$} & {$[12; 15)$} & {$[15; 18)$} & {$[18; 21)$} & {$[21; 24)$} \\
\hline Số nhân viên & $6$ & $12$ & $4$ & $2$ & $1$ \\
\hline
\end{tabular}
\end{center}
Xét tính đúng sai của các mệnh đề sau:
	\choiceTF
	{\True Giá trị đại diện của nhóm $[9; 12)$ là $10{,}5$}
	{Trung bình lương các nhân viên là $16{,}5$ triệu đồng}
	{Nhóm chứa trung vị là $[15; 18)$}
	{\True Tứ phân vị thứ ba là $15{,}56$}
	\loigiai{
	\begin{itemchoice}
	\itemch Đúng. Vì giá trị đại diện của nhóm $[9; 12)$ là $\dfrac{9+12}{2}=10{,}5$.
	\itemch Sai. Vì trung bình lương các nhân viên là $$\overline{x}=\dfrac{1}{25}(6 \cdot 10{,}5+12 \cdot 13{,}5+4\cdot 16{,}5+2\cdot 19{,}5+22{,}5)=14{,}1\text{ triệu đồng.}$$ 
	\itemch Sai. Vì công ty có $25$ nhân sự. Vì $x_{13}\in[12; 15)$ nên nhóm này chứa trung vị. 
	\itemch Đúng. Vì $x_{19}\in[15; 18)$ nên ta có $Q_3=15+\dfrac{\dfrac{3 \cdot 25}{4}-18}{4} \cdot 3\approx 15{,}56$.
	\end{itemchoice}
}
\end{ex}

\begin{ex}%[1D5V2-3]%Câu 3
Điều tra về chiều cao của học sinh khối lớp 10, ta có kết quả sau:
\begin{center}
\begin{tabular}{|c|c|c|}
\hline Nhóm & Chiều cao $(\mathrm{cm})$ & Số học sinh \\
\hline 1 & {$[150; 152)$} & $5$ \\
\hline 2 & {$[152; 154)$} & $18$ \\
\hline 3 & {$[154; 156)$} & $40$ \\
\hline 4 & {$[156; 158)$} & $26$ \\
\hline 5 & {$[158; 160)$} & $8$ \\
\hline 6 & {$[160; 162)$} & $3$ \\
\hline
\end{tabular}
\end{center}
	\choiceTF
	{\True Tổng số học sinh điều tra bằng $100$}
	{\True Chiều cao trung bình của các em học sinh trên là $155{,}46$}
	{Số trung vị của mẫu số liệu trên bằng $156{,}35$}
	{\True Tứ phân vị thứ nhất của mẫu số liệu trên bằng $154{,}1$}
	\loigiai{
	\begin{itemchoice}
	\itemch Đúng. Vì ta có tổng số học sinh là $n=100$.
	\itemch Đúng. Vì chiều cao trung bình bằng $$\overline{x}=\dfrac{5\cdot 151+18\cdot 153+40\cdot 155+26\cdot 157+8.159+3\cdot 161}{100}=155{,}46.$$ 
	\itemch Sai. Vì $n=100$ nên trung vị là $\dfrac{x_{50}+x_{51}}{2}$, do hai giá trị $x_{50}, x_{51}$ thuộc nhóm 3.\\
Áp dụng công thức $M_e=154+\dfrac{50-(18+5)}{40} \cdot 2 \approx 155{,}35$.
	\itemch Đúng. Vì tứ phân vị thứ nhất $Q_3$ là $\dfrac{x_{25}+x_{26}}{2}$ mà $x_{25}, x_{26}$ thuộc nhóm $[154; 156)$.\\
Do đó, $p=3$, $a_p=154$, $m_p=40$, $m_1+m_2=5+18=23$, $a_{p+1}-a_p=2$ và ta có
$$Q_1=154+\dfrac{\dfrac{100}{4}-23}{40} \times 2=154{,}1.$$ 
	\end{itemchoice}
}
\end{ex}

\begin{ex}%[1D5V2-3]%Câu 4
Kết quả điều tra về số giờ làm thêm trong một tuần của sinh viên một trường đại học $\mathrm{X}$ được cho bởi bảng sau:
\begin{center}
\begin{tabular}{|l|c|c|c|c|c|}
\hline Số giờ làm thêm & {$[2; 4)$} & {$[4; 6)$} & {$[6; 8)$} & {$[8; 10)$} & {$[10; 12)$}\\
\hline Số sinh viên & $12$ & $20$ & $37$ & $21$ & $10$ \\
\hline
\end{tabular}
\end{center}
	\choiceTF
	{\True Số sinh viên được điều tra là $100$}
	{\True Số giờ làm thêm trung bình của mỗi sinh viên trường đại học $\mathrm{X}$ không ít hơn $6$}
	{Mốt của mẫu số liệu trên là $7{,}5$}
	{Tứ phân vị thứ hai của dãy số liệu lớn hơn $6{,}5$}
	\loigiai{
	\begin{itemchoice}
	\itemch Đúng. Vì cỡ mẫu $n=100$.
	\itemch Đúng. Vì từ bảng thống kê ta có
\begin{center}
\begin{tabular}{|l|c|c|c|c|c|}
\hline Số giờ làm thêm & {$[2; 4)$} & {$[4; 6)$} & {$[6; 8)$} & {$[8; 10)$} & {$[10; 12)$} \\
\hline Số giờ làm thêm đai diện & $3$ & $5$ & $7$ & $9$ & $11$ \\
\hline Số sinh viên & $12$ & $20$ & $37$ & $21$ & $10$ \\
\hline
\end{tabular}
\end{center}
Số trung bình của mẫu số liệu là
$$\overline{x}=\dfrac{3\cdot 12+5\cdot 20+7\cdot 37+9\cdot 21+11\cdot 10}{100}=6\cdot 94.$$
	\itemch Sai. Vì nhóm chứa mốt số liệu trên là nhóm $[6; 8)$.\\
Do đó $u_m=6$; $n_m=37$, $n_{m-1}=20$; $n_{m+1}=21$; $u_{m+1}=8$.\\
Vậy mốt của mẫu số liệu ghép nhóm là
$$M_0=6+\dfrac{37-20}{(37-20)+(37-21)}\cdot(8-6)\approx 7{,}03.$$ 
	\itemch Sai. Vì $x_1; x_2; \ldots; x_{100}$ là mẫu số liệu theo thứ tự không giảm.\\
Tứ phân vị thứ hai của dãy số liệu $x_1; x_2; \ldots; x_{100}$ là $\dfrac{1}{2}\left(x_{50}+x_{51}\right)$.\\
Do $x_{50}$ và $x_{51}$ thuộc nhóm $[6; 8)$ nên $Q_2=6+\dfrac{\dfrac{2 \cdot 100}{4}-32}{37} \cdot(8-6) \approx 6{,}97$.
	\end{itemchoice}
}
\end{ex}
\Closesolutionfile{ans}
\indapan{3}{ans/ans-OTC3-DS}

\Opensolutionfile{ans}[ans/ans-OTC3-KQ]
\caukq
\setcounter{ex}{0}
\begin{ex}%[1D5H1-3]%Câu 1
Khi đo cân nặng của học sinh lớp 11D, y tá lập được bảng số liệu ghép nhóm sau đây:
\begin{center}
\begin{tabular}{|l|c|c|c|c|c|c|}
\hline Cân nặng (kg) & {$[40{,}5; 45{,}5)$} & {$[45{,}5; 50{,}5)$} & {$[50{,}5; 55{,}5)$} & {$[55{,}5; 60{,}5)$} & {$[60{,}5; 65{,}5)$} & {$[65{,}5; 70{,}5)$} \\
\hline Số học sinh & $10$ & $7$ & $16$ & $4$ & $2$ & $3$ \\
\hline
\end{tabular}
\end{center}
Cân nặng trung bình của học sinh lớp 11D là bao nhiêu? (Kết quả làm tròn đến hàng phần chục)
\shortans{$51{,}8$}
	\loigiai{
Ta có bảng thống kê
\begin{center}
\begin{tabular}{|l|c|c|c|c|c|c|}
\hline Cân nặng $(\mathrm{kg})$ & {$[40{,}5; 45{,}5)$} & {$[45{,}5; 50{,}5)$} & {$[50{,}5 ; 55{,}5)$} & {$[55{,}5; 60{,}5)$} & {$[60{,}5; 65{,}5)$} & {$[65{,}5 ; 70{,}5)$} \\
\hline Giá trị đại diện & $43$ & $48$ & $53$ & $58$ & $63$ & $68$ \\
\hline Số học sinh & $10$ & $7$ & $16$ & $4$ & $2$ & $3$ \\
\hline
\end{tabular}
\end{center}
Cân nặng trung bình của học sinh là $$\overline{x}=\dfrac{43\cdot 10+48\cdot 7+53\cdot 16+58\cdot 4+63\cdot 2+68\cdot 3}{42}=51{,}8.$$
}
\end{ex}

\begin{ex}%[1D5V1-3]%Câu 2
Thời gian truy cập internet mỗi buổi trưa của một số học sinh được cho trong bảng sau:
\begin{center}
\begin{tabular}{|l|c|c|c|c|c|}
\hline Thời gian (phút) & {$[9{,}5; 12{,}5)$} & {$[12{,}5; 15{,}5)$} & {$[15{,}5 ; 18{,}5)$} & {$[18{,}5; 21{,}5)$} & {$[21{,}5; 24{,}5)$} \\
\hline Số học sinh & $3$ & $12$ & $15$ & $24$ & $2$ \\
\hline
\end{tabular}
\end{center}
Số trung vị của mẫu số liệu trên là bao nhiêu? (Kết quả làm tròn đến hàng phần chục)
\shortans{$18{,}1$}
	\loigiai{
Cỡ mẫu là $n=3+12+15+24+2=56$.\\
Gọi $x_1; \ldots; x_{56}$ là thời gian truy cập internet của $56$ học sinh và giả sử dãy này đã được sắp xếp theo thứ tự tăng dần.\\
Khi đó, số trung vị là $\dfrac{x_{28}+x_{29}}{2}$.\\
Do hai giá trị $x_{28}; x_{29}$ thuộc nhóm $[15{,}5; 18{,}5)$ nên nhóm này chứa trung vị.\\
Do đó $p=3$; $a_3=15{,}5$; $m_3=15$; $m_1+m_2=15$; $a_4-a_3=3$.
$$M_e=15{,}5+\dfrac{\dfrac{56}{2}-15}{15}\cdot 3=18{,}1.$$
}
\end{ex}

\begin{ex}%[1D5V2-3]%Câu 3
Khảo sát thời gian tập thể dục của một số học sinh khối 11 thu được mẫu số liệu ghép nhóm sau:
\begin{center}
\begin{tabular}{|l|c|c|c|c|c|}
\hline Thời gian (phút) & {$[0; 20)$} & {$[20; 40)$} & {$[40; 60)$} & {$[60; 80)$} & {$[80; 100)$} \\
\hline Số học sinh & $5$ & $9$ & $12$ & $10$ & $6$ \\
\hline
\end{tabular}
\end{center}
Tứ phân vị thứ ba của mẫu số liệu là bao nhiêu?
\shortans{$71$}
	\loigiai{
Cỡ mẫu là $n=5+9+12+10+6=42$.\\
Tứ phân vị thứ ba $Q_3$ là $x_{33}$, do nhóm chứa $x_{33}$ là $[60; 80)$, do đó $p=4$; $a_4=60$; $m_1+m_2+m_3=5+9+12=26$; $a_5-a_4=20$.\\
Ta có $Q_3=a_4+\dfrac{\dfrac{3 n}{4}-\left(m_1+m_2+m_3\right)}{m_4}\left(a_5-a_4\right)=60+\dfrac{\dfrac{3.42}{4}-26}{10} \cdot 20=71$.
}
\end{ex}

\begin{ex}%[1D5H1-3]%Câu 4
Khảo sát chiều cao của học sinh lớp $11 \mathrm{C} 1$, thu được bảng số liệu ghép nhóm như sau
\begin{center}
\begin{tabular}{|l|c|c|c|c|c|}
\hline Khoàng chiều cao $(\mathrm{cm})$ & {$[150; 155)$} & {$[155; 160)$} & {$[160; 165)$} & {$[165; 170)$} & {$[170; 175)$} \\
\hline Số học sinh & $7$ & $15$ & $12$ & $8$ & $3$ \\
\hline
\end{tabular}
\end{center}
Chiều cao trung bình của học sinh 11C1 có dạng $abc{,}d$, với $a,b,c,d$ là các số tự nhiên. Tính $S=a+b+c+d$.
\shortans{$15$}
	\loigiai{
\begin{center}
\begin{tabular}{|l|c|c|c|c|c|}
\hline Khoàng chiều cao $(\mathrm{cm})$ & {$[150; 155)$} & {$[155; 160)$} & {$[160; 165)$} & {$[165; 170)$} & {$[170; 175)$} \\
\hline Giá trị đại diện & $152{,}5$ & $157{,}5$ & $162{,}5$ & $167{,}5$ & $172{,}5$ \\
\hline Số học sinh & $7$ & $15$ & $12$ & $8$ & $3$ \\
\hline
\end{tabular}
\end{center}
Chiều cao trung bình của học sinh lớp 11C1 là
$$\overline{x}=\dfrac{1}{45}(7\cdot 152{,}5+15\cdot 157{,}5+12\cdot 162{,}5+4 \cdot 167{,}5+3\cdot 172{,}5)=160{,}8 \mathrm{~cm}.$$
Khi đó $S=1+6+0+8=15$.
}
\end{ex}

\begin{ex}%[1D5H1-3]%Câu 5
Doanh thu bán hàng trong 20 ngày được lựa chọn ngẫu nhiên của một cửa hàng được ghi lại ở bảng sau (đơn vị: triệu đồng):
\begin{center}
\begin{tabular}{|l|c|c|c|c|c|}
\hline Doanh thu & {$[5; 7)$} & {$[7; 9)$} & {$[9; 11)$} & {$[11; 13)$} & {$[13; 15)$} \\
\hline Số ngày & $2$ & $7$ & $7$ & $3$ & $1$ \\
\hline
\end{tabular}
\end{center}
Doanh số trung bình của cửa hàng trên trong $20$ ngày là bao nhiêu? (Kết quả làm tròn đến hàng phần chục)
\shortans{$9{,}4$}
	\loigiai{
Bảng tần số ghép nhóm theo giá trị đại diện là
\begin{center}
\begin{tabular}{|l|c|c|c|c|c|}
\hline Doanh thu & {$[5; 7)$} & {$[7; 9)$} & {$[9; 11)$} & {$[11; 13)$} & {$[13; 15)$} \\
\hline Giá trị đại diện & $6$ & $8$ & $10$ & $12$ & $14$ \\
\hline Số ngày & $2$ & $7$ & $7$ & $3$ & $1$ \\
\hline
\end{tabular}
\end{center}
Số trung bình $\overline{x}=\dfrac{2\cdot 6+7\cdot 8+7\cdot 10+3\cdot 12+1\cdot 14}{20}=9{,}4$.
}
\end{ex}

\begin{ex}%[1D5V1-4]%Câu 6
Một công ty bất động sản Đất Vàng thực hiện cuộc khảo sát khách hàng xem họ có nhu cầu mua nhà ở mức giá nào để tiến hành dự án xây nhà ở Thăng Long group sắp tới. Kết quả khảo sát $500$ khách hàng được ghi lại ở bảng sau:
\begin{center}
\begin{tabular}{|l|c|c|c|c|c|}
\hline Mức giá (triệu đồng) & {$[10; 14)$} & {$[14; 18)$} & {$[18; 22)$} & {$[22; 26)$} & {$[26; 30)$} \\
\hline Số khách hàng & $75$ & $105$ & $179$ & $96$ & $45$ \\
\hline
\end{tabular}
\end{center}
Công ty bất động sản Đất Vàng nên xây nhà ở mức giá nào để nhiều người có nhu cầu xây nhà?
\shortans{$19{,}9$}
	\loigiai{
Nhóm chứa mốt của mẫu số liệu trên là $[18; 22)$.\\
Do đó $u_m=18$; $n_{m-1}=105$; $n_{m+1}=96$; $u_{m+1}-u_m=22-18=4$.\\
Mốt của mẫu số liệu ghép nhóm là $M_0=18+\dfrac{179-105}{(179-105)+(179-96)} \cdot 4 \approx 19{,}9$.\\
Dựa vào kết quả trên ta có thể dự đoán rằng nếu công ty bất động sản Đất Vàng xây nhà ở mức giá $19{,}9$ triệu đồng$/\mathrm{m}^2$ thì sẽ có nhiều người mua nhất.
}
\end{ex}
\Closesolutionfile{ans}
\indapan{6}{ans/ans-OTC3-KQ}

\Opensolutionfile{ans}[ans/ans-OTC3-TN]
\caulc
\begin{ex}%[1D5N1-2]%Câu 1
Điều tra về chiều cao của học sinh khối lớp 11, ta được mẫu số liệu sau:
\begin{center}
\begin{tabular}{|c|c|}
\hline Chiều cao $(\mathrm{cm})$ & Số học sinh \\
\hline$[150; 152)$ & $5$ \\
\hline$[152; 154)$ & $18$ \\
\hline$[154; 156)$ & $40$ \\
\hline$[156; 158)$ & $26$ \\
\hline$[158; 160)$ & $8$ \\
\hline$[160; 162)$ & $3$ \\
\hline Tổng & $N=100$ \\
\hline
\end{tabular}
\end{center}
Mẫu số liệu ghép nhóm đã cho có tất cả bao nhiêu nhóm?
	\choice
	{$5$}
	{\True $6$}
	{$7$}
	{$12$}
	\loigiai{
Mẫu số liệu ghép nhóm đã cho có tất cả $6$ nhóm.
}
\end{ex}

\begin{ex}%[1D5N1-2]%Câu 2
Cho mẫu số liệu về chiều cao $(\mathrm{cm})$ của các học $\sinh$ nữ trong khối 11 của một trường như sau:
\begin{center}
\begin{tabular}{|l|c|c|c|c|c|c|}
\hline Chiều cao & {$[145; 150)$} & {$[150; 155)$} & {$[155; 160)$} & {$[160; 165)$} & {$[165; 170)$} & {$[170; 175)$} \\
\hline Số học sinh & $20$ & $45$ & $34$ & $27$ & $15$ & $4$ \\
\hline
\end{tabular}
\end{center}
Mẫu số liệu trên có bao nhiêu số liệu, bao nhiêu nhóm?
	\choice
	{$145$ số liệu; $6$ nhóm}
	{$30$ số liệu; $5$ nhóm}
	{\True $6$ số liệu; $145$ nhóm}
	{$5$ số liệu; $30$ nhóm}
	\loigiai{
Mẫu số liệu $(T)$ có
$$20+45+34+27+15+4=145\text{ (số liệu).}$$
$6$ nhóm là $[145; 150)$; $[150; 155)$; $[155; 160)$; $[160; 165)$; $[165; 170)$; $[170; 175)$.
}
\end{ex}

\begin{ex}%[1D5N1-2]%Câu 3
Tìm hiểu thời gia xem tivi trong tuần trước (đơn vị: giờ) của một số học sinh thu được kết quả sau:
\begin{center}
\begin{tabular}{|l|c|c|c|c|c|}
\hline Thời gian (giờ) & {$[0; 5)$} & {$[5; 10)$} & {$[10; 15)$} & {$[15; 20)$} & {$[20; 25)$} \\
\hline Số học sinh & $8$ & $16$ & $4$ & $2$ & $2$ \\
\hline
\end{tabular}
\end{center}
Giá trị đại diện của nhóm $[20; 25)$ là
	\choice
	{\True $22{,}5$}
	{$23$}
	{$20$}
	{$5$}
	\loigiai{
Giá trị đại diện của nhóm $[20; 25)$ là $\dfrac{20+25}{2}=22{,}5$.
}
\end{ex}

\begin{ex}%[1D5N1-4]%Câu 4
Khảo sát thời gian tập thể dục của một số học sinh khối 11 thu được mẫu số liệu ghép nhóm sau:
\begin{center}
\begin{tabular}{|l|c|c|c|c|c|}
\hline Thời gian (phút) & {$[0; 20)$} & {$[20; 40)$} & {$[40; 60)$} & {$[60; 80)$} & {$[80; 100)$} \\
\hline Số học sinh & $5$ & $9$ & $12$ & $10$ & $6$\\
\hline
\end{tabular}
\end{center}
Nhóm chứa mốt của mẫu số liệu trên là
	\choice
	{\True $[40; 60)$}
	{$[20; 40)$}
	{$[60; 80)$}
	{$[80; 100)$}
	\loigiai{
Mốt $M_0$ chứa trong nhóm $[40; 60)$.
}
\end{ex}

\begin{ex}%[1D5H2-2]%Câu 5
Khảo sát thời gian tập thể dục của một số học sinh khối 11 thu được mẫu số liệu ghép nhóm sau:
\begin{center}
\begin{tabular}{|l|c|c|c|c|c|}
\hline Thời gian (phút) & {$[0; 20)$} & {$[20; 40)$} & {$[40; 60)$} & {$[60; 80)$} & {$[80; 100)$} \\
\hline Số học sinh & $5$ & $9$ & $12$ & $10$ & $6$ \\
\hline
\end{tabular}
\end{center}
Nhóm chứa trung vị của mẫu số liệu trên là
	\choice
	{\True $[40; 60)$}
	{$[20; 40)$}
	{$[60; 80)$}
	{$[80; 100)$}
	\loigiai{
Ta có $n=42$.\\
Nên trung vị của mẫu số liệu trên là $Q_2=\dfrac{x_{21}+x_{22}}{2}$.\\
Mà $x_{21}, x_{22} \in[40; 60)$ nên nhóm chứa trung vị của mẫu số liệu trên là nhóm $[40; 60)$.
}
\end{ex}

\begin{ex}%[1D5H2-3]%Câu 6
Khảo sát thời gian tập thể dục của một số học sinh khối 11 thu được mẫu số liệu ghép nhóm sau:
\begin{center}
\begin{tabular}{|l|c|c|c|c|c|}
\hline Thời gian (phút) & {$[0; 20)$} & {$[20; 40)$} & {$[40; 60)$} & {$[60; 80)$} & {$[80; 100)$} \\
\hline Số học sinh & $5$ & $9$ & $12$ & $10$ & $6$ \\
\hline
\end{tabular}
\end{center}
Nhóm chứa tứ phân vị thứ ba của mẫu số liệu trên là
	\choice
	{$[40; 60)$}
	{$[20; 40)$}
	{\True $[60; 80)$}
	{$[80; 100)$}
	\loigiai{
Ta có $n=42$.\\
Nên tứ phân vị thứ nhất của mẫu số liệu trên là $Q_3=x_{33}$.\\
Mà $x_{33} \in[60; 80)$ nên nhóm chứa tứ phân vị thứ nhất của mẫu số liệu trên là nhóm $[60; 80)$.
}
\end{ex}

\begin{ex}%[1D5H1-3]%Câu 7
Doanh thu bán hàng trong $20$ ngày được lựa chọn ngẫu nhiên của một của hàng được ghi lại ở bảng sau (đơn vị: triệu đồng).
\begin{center}
\begin{tabular}{|l|c|c|c|c|c|}
\hline Doanh thu & {$[5; 7)$} & {$[7; 9)$} & {$[9; 11)$} & {$[11; 13)$} & {$[13; 15)$} \\
\hline Số ngày & $2$ & $7$ & $7$ & $3$ & $1$ \\
\hline
\end{tabular}
\end{center}
Số trung bình của mẫu số liệu trên thuộc khoảng nào trong các khoảng dưới đây?
	\choice
	{$[7; 9)$}
	{\True $[9; 11)$}
	{$[11; 13)$}
	{$[13; 15)$}
	\loigiai{
Bảng tần số ghép nhóm theo giá trị đại diện là
\begin{center}
\begin{tabular}{|l|c|c|c|c|c|}
\hline Doanh thu & {$[5; 7)$} & {$[7; 9)$} & {$[9; 11)$} & {$[11; 13)$} & {$[13; 15)$} \\
\hline Giá trị đại diện & $6$ & $8$ & $10$ & $12$ & $14$ \\
\hline Số ngày & $2$ & $7$ & $7$ & $3$ & $1$ \\
\hline
\end{tabular}
\end{center}
Số trung bình là $\overline{x}=\dfrac{2\cdot6+7\cdot8+7\cdot10+3\cdot12+1\cdot14}{20}=9{,}4$.
}
\end{ex}

\begin{ex}%[1D5V2-2]%Câu 8
Cho mẫu số liệu ghép nhóm về thống kê điểm số (thang điểm $20$) của 100 học sinh tham dự kỳ thi học sinh giỏi toán, ta có bảng số liệu sau:
\begin{center}
\begin{tabular}{|l|c|c|c|c|c|c|}
\hline Điểm & {$[8; 10)$} & {$[10; 12)$} & {$[12; 14)$} & {$[14; 16)$} & {$[16; 18)$} & {$[18; 20)$} \\
\hline Số học sinh & $6$ & $21$ & $30$ & $25$ & $14$ & $4$\\
\hline
\end{tabular}
\end{center}
Tìm trung vị của mẫu số liệu ghép nhóm trên. (Kết quả làm tròn đến hàng phần trăm).
	\choice
	{$12{,}18$}
	{$12{,}81$}
	{$13{,}35$}
	{\True $13{,}53$}
	\loigiai{
Cỡ mẫu là $n=6+21+30+25+14+4=100$.\\
Gọi $x_1; x_2; \ldots; x_{100}$ là số điểm của 100 học sinh và được sắp xếp theo thứ tự tăng dần. Khi đó, trung vị là $\dfrac{x_{50}+x_{51}}{2}$, do hai giá trị $x_{50}, x_{51}$ thuộc nhóm $[12; 14)$ nên nhóm này chứa trung vị.\\
Ta có $p=3$; $a_3=12$; $m_3=30$; $m_1+m_2=25$; $a_4-a_3=2$ khi đó trung vị của mẫu số liệu là
$$M_e=12+\dfrac{\dfrac{100}{2}-27}{30}\cdot 2=13{,}53.$$
}
\end{ex}

\begin{ex}%[1D5N1-4]%Câu 9
Cho mẫu số liệu về thời gian (phút) đi từ nhà đến trường của các học sinh trong một lớp 11 của một trường như sau:
\begin{center}
\begin{tabular}{|l|c|c|c|c|c|c|}
\hline Thời gian & {$[0; 5)$} & {$[5; 10)$} & {$[10; 15)$} & {$[15; 20)$} & {$[20; 25)$} & {$[25; 30)$} \\
\hline Số học sinh & $7$ & $12$ & $7$ & $5$ & $3$ & $2$ \\
\hline
\end{tabular}
\end{center}
Có bao nhiêu học sinh có thời gian đi từ nhà đến trường là $15$ phút đến $20$ phút?
	\choice
	{$20$}
	{$15$}
	{\True $5$}
	{$7$}
	\loigiai{
Theo bảng số liệu, có $5$ học sinh có thời gian đi từ nhà đến trường là $15$ phút đến $20$ phút.
}
\end{ex}

\begin{ex}%[1D5H1-3]%Câu 10
Cho mẫu số liệu ghép nhóm về thống kê nhiệt độ tại một địa điểm trong $30$ ngày, ta có bảng số liệu sau:
\begin{center}
\begin{tabular}{|l|c|c|c|c|}
\hline Nhiệt độ $\left({ }^{\circ} \mathrm{C}\right)$ & {$[18; 21)$} & {$[21; 24)$} & {$[24; 27)$} & {$[27; 30)$}\\
\hline Số ngày & $6$ & $12$ & $9$ & $3$ \\
\hline
\end{tabular}
\end{center}
Nhiệt độ trung bình trong $30$ ngày trên là
	\choice
	{$24^{\circ} \mathrm{C}$}
	{$25{,}4^{\circ} \mathrm{C}$}
	{$24{,}3^{\circ} \mathrm{C}$}
	{\True $23{,}4^{\circ} \mathrm{C}$}
	\loigiai{
Giá trị đại diện của mỗi nhóm số liệu là trung bình cộng của hai đầu mút.\\
Ta có bảng tần số ghép nhóm theo giá trị đại diện của mỗi nhóm:
\begin{center}
\begin{tabular}{|l|c|c|c|c|}
\hline Nhóm & {$[18; 21)$} & {$[21; 24)$} & {$[24; 27)$} & {$[27; 30)$}\\
\hline Giá trị đại diện & $19{,}5$ & $22{,}5$ & $25{,}5$ & $28{,}5$ \\
\hline Tần số & $6$ & $12$ & $9$ & $3$ \\
\hline
\end{tabular}
\end{center}
Nhiệt độ trung bình trong 30 ngày trên là
$$\overline{x}=\dfrac{6 \cdot 19{,}5+12 \cdot 22{,}5+9\cdot25{,}5+3 \cdot 28{,}5}{30}=23{,}4^{\circ} \mathrm{C}.$$
}
\end{ex}

\begin{ex}%[1D5N1-4]%Câu 11
Cho mẫu số liệu về chiều cao $(\mathrm{cm})$ của các học sinh nữ trong khối 11 của một trường như sau:
\begin{center}
\begin{tabular}{|l|c|c|c|c|c|c|}
\hline Chiều cao & {$[145; 150)$} & {$[150; 155)$} & {$[155; 160)$} & {$[160; 165)$} & {$[165; 170)$} & {$[170; 175)$} \\
\hline Số học sinh & $20$ & $45$ & $34$ & $27$ & $15$ & $4$ \\
\hline
\end{tabular}
\end{center}
Số học sinh nữ cao từ $150 \mathrm{~cm}$ đến $155 \mathrm{~cm}$ là
	\choice
	{$20$}
	{$65$}
	{$34$}
	{\True $45$}
	\loigiai{
Theo bảng số liệu trên, số học sinh nữ cao từ $150 \mathrm{~cm}$ đến $155 \mathrm{~cm}$ là $45$ học sinh.
}
\end{ex}

\begin{ex}%[1D5V2-2]%Câu 12
Cho mẫu số liệu ghép nhóm về chiều cao của $25$ cây dừa giống như sau:
\begin{center}
\begin{tabular}{|l|c|c|c|c|c|}
\hline Chiều cao $(\mathrm{cm})$ & {$[0; 10)$} & {$[10; 20)$} & {$[20; 30)$} & {$[30; 40)$} & {$[40; 50)$} \\
\hline Số cây & $4$ & $6$ & $7$ & $5$ & $3$ \\
\hline
\end{tabular}
\end{center}
Trung vị của mẫu số liệu ghép nhóm này là
	\choice
	{$M_e=\dfrac{175}{7}$}
	{$M_e=\dfrac{165}{5}$}
	{$M_e=\dfrac{165}{7}$}
	{\True $M_e=\dfrac{165}{3}$}
	\loigiai{
Cỡ mẫu $n=4+6+7+5+3=25$.\\
$x_1, x_2, \ldots, x_{25}$ là chiều cao của $25$ cây dừa giống được sắp xếp theo thứ tự không giảm. Khi đó, trung vị là $x_{13}$. Do $x_{13}$ thuộc nhóm $[20; 30)$ nên nhóm này chứa trung vị.\\
 Do đó $p=3$, $a_3=20$, $m_3=7$, $m_1+m_2=10$, $a_4-a_3=10$.\\
Vậy $M_e=20+\dfrac{\dfrac{25}{2}-10}{7} \cdot 10=\dfrac{165}{7}$.
}
\end{ex}
\Closesolutionfile{ans}
\indapan{6}{ans/ans-OTC3-TN}

\cauds
\Opensolutionfile{ans}[ans/ans-OTC3-DS]
\setcounter{ex}{0}
\begin{ex}%[1D5H1-4]%Câu 1
Cho mẫu số liệu điểm môn Toán của một nhóm học sinh như sau:
\begin{center}
\begin{tabular}{|l|c|c|c|c|}
\hline Điểm & $[6; 7)$ & $[7; 8)$ & $[8; 9)$ & $[9; 10]$ \\
\hline Số học sinh & $8$ & $7$ & $10$ & $5$ \\
\hline
\end{tabular}
\end{center}
Xét tính đúng sai của các mệnh đề sau:
	\choiceTF
	{\True Mẫu số liệu đã cho là mẫu số liệu ghép nhóm}
	{\True Cỡ mẫu của mẫu số liệu là $30$}
	{\True Điểm trung bình của các học sinh là $7{,}9$}
	{Mốt của mẫu số liệu là $10$}
	\loigiai{
Ta có bảng sau
\begin{center}
\begin{tabular}{|l|c|c|c|c|}
\hline Điểm & {$[6; 7)$} & {$[7; 8)$} & {$[8; 9)$} & {$[9; 10]$} \\
\hline Giá trị đại diện & $6{,}5$ & $7{,}5$ & $8{,}5$ & $9{,}5$ \\
\hline Số học sinh & $8$ & $7$ & $10$ & $5$ \\
\hline
\end{tabular}
\end{center}
	\begin{itemchoice}
	\itemch Đúng. Vì theo định nghĩa mẫu số liệu ghép nhóm.
	\itemch Đúng. Vì cỡ mẫu là $8+7+10+5=30$.
	\itemch Đúng. Vì điểm trung bình của các học sinh là $$\overline{x}=\dfrac{6{,}58+7{,}5 \cdot 7+8{,}5 \cdot 10+9{,}5 \cdot 5}{30}=7{,}9.$$
	\itemch Sai. Vì nhóm chứa Mốt là $[8; 9)$.\\
 Mốt của mẫu số liệu là $M_e=8+\dfrac{10-7}{10-7+10-5}(9-8)=8{,}375$.
	\end{itemchoice}
}
\end{ex}

\begin{ex}%[1D5H2-3]%Câu 2
Cho mẫu số liệu ghép nhóm về lương của nhân viên trong một công ty như sau:
\begin{center}
\begin{tabular}{|l|c|c|c|c|c|}
\hline Lương (triệu đồng) & {$[9; 12)$} & {$[12; 15)$} & {$[15; 18)$} & {$[18; 21)$} & {$[21; 24)$} \\
\hline Số nhân viên & $6$ & $12$ & $4$ & $2$ & $1$ \\
\hline
\end{tabular}
\end{center}
Xét tính đúng sai của các mệnh đề sau:
	\choiceTF
	{\True Giá trị đại diện của nhóm $[9; 12)$ là $10{,}5$}
	{Trung bình lương các nhân viên là $16{,}5$ triệu đồng}
	{Nhóm chứa trung vị là $[15; 18)$}
	{\True Tứ phân vị thứ ba là $15{,}56$}
	\loigiai{
	\begin{itemchoice}
	\itemch Đúng. Vì giá trị đại diện của nhóm $[9; 12)$ là $\dfrac{9+12}{2}=10{,}5$.
	\itemch Sai. Vì trung bình lương các nhân viên là $$\overline{x}=\dfrac{1}{25}(6 \cdot 10{,}5+12 \cdot 13{,}5+4\cdot 16{,}5+2\cdot 19{,}5+22{,}5)=14{,}1\text{ triệu đồng.}$$ 
	\itemch Sai. Vì công ty có $25$ nhân sự. Vì $x_{13}\in[12; 15)$ nên nhóm này chứa trung vị. 
	\itemch Đúng. Vì $x_{19}\in[15; 18)$ nên ta có $Q_3=15+\dfrac{\dfrac{3 \cdot 25}{4}-18}{4} \cdot 3\approx 15{,}56$.
	\end{itemchoice}
}
\end{ex}

\begin{ex}%[1D5V2-3]%Câu 3
Điều tra về chiều cao của học sinh khối lớp 10, ta có kết quả sau:
\begin{center}
\begin{tabular}{|c|c|c|}
\hline Nhóm & Chiều cao $(\mathrm{cm})$ & Số học sinh \\
\hline 1 & {$[150; 152)$} & $5$ \\
\hline 2 & {$[152; 154)$} & $18$ \\
\hline 3 & {$[154; 156)$} & $40$ \\
\hline 4 & {$[156; 158)$} & $26$ \\
\hline 5 & {$[158; 160)$} & $8$ \\
\hline 6 & {$[160; 162)$} & $3$ \\
\hline
\end{tabular}
\end{center}
	\choiceTF
	{\True Tổng số học sinh điều tra bằng $100$}
	{\True Chiều cao trung bình của các em học sinh trên là $155{,}46$}
	{Số trung vị của mẫu số liệu trên bằng $156{,}35$}
	{\True Tứ phân vị thứ nhất của mẫu số liệu trên bằng $154{,}1$}
	\loigiai{
	\begin{itemchoice}
	\itemch Đúng. Vì ta có tổng số học sinh là $n=100$.
	\itemch Đúng. Vì chiều cao trung bình bằng $$\overline{x}=\dfrac{5\cdot 151+18\cdot 153+40\cdot 155+26\cdot 157+8.159+3.161}{100}=155{,}46.$$ 
	\itemch Sai. Vì $n=100$ nên trung vị là $\dfrac{x_{50}+x_{51}}{2}$, do hai giá trị $x_{50}, x_{51}$ thuộc nhóm 3.\\
Áp dụng công thức $M_e=154+\dfrac{50-(18+5)}{40} \cdot 2 \approx 155{,}35$.
	\itemch Đúng. Vì tứ phân vị thứ nhất $Q_3$ là $\dfrac{x_{25}+x_{26}}{2}$ mà $x_{25}, x_{26}$ thuộc nhóm $[154; 156)$.\\
Do đó, $p=3$, $a_p=154$, $m_p=40$, $m_1+m_2=5+18=23$, $a_{p+1}-a_p=2$ và ta có
$$Q_1=154+\dfrac{\dfrac{100}{4}-23}{40} \times 2=154{,}1.$$ 
	\end{itemchoice}
}
\end{ex}

\begin{ex}%[1D5V2-3]%Câu 4
Kết quả điều tra về số giờ làm thêm trong một tuần của sinh viên một trường đại học $\mathrm{X}$ được cho bởi bảng sau:
\begin{center}
\begin{tabular}{|l|c|c|c|c|c|}
\hline Số giờ làm thêm & {$[2; 4)$} & {$[4; 6)$} & {$[6; 8)$} & {$[8; 10)$} & {$[10; 12)$}\\
\hline Số sinh viên & $12$ & $20$ & $37$ & $21$ & $10$ \\
\hline
\end{tabular}
\end{center}
	\choiceTF
	{\True Số sinh viên được điều tra là $100$}
	{\True Số giờ làm thêm trung bình của mỗi sinh viên trường đại học $\mathrm{X}$ không ít hơn $6$}
	{Mốt của mẫu số liệu trên là $7{,}5$}
	{Tứ phân vị thứ hai của dãy số liệu lớn hơn $6{,}5$}
	\loigiai{
	\begin{itemchoice}
	\itemch Đúng. Vì cỡ mẫu $n=100$.
	\itemch Đúng. Vì từ bảng thống kê ta có
\begin{center}
\begin{tabular}{|l|c|c|c|c|c|}
\hline Số giờ làm thêm & {$[2; 4)$} & {$[4; 6)$} & {$[6; 8)$} & {$[8; 10)$} & {$[10; 12)$} \\
\hline Số giờ làm thêm đai diện & $3$ & $5$ & $7$ & $9$ & $11$ \\
\hline Số sinh viên & $12$ & $20$ & $37$ & $21$ & $10$ \\
\hline
\end{tabular}
\end{center}
Số trung bình của mẫu số liệu là
$$\overline{x}=\dfrac{3\cdot 12+5\cdot 20+7\cdot 37+9\cdot 21+11\cdot 10}{100}=6\cdot 94.$$
	\itemch Sai. Vì nhóm chứa mốt số liệu trên là nhóm $[6; 8)$.\\
Do đó $u_m=6$; $n_m=37$, $n_{m-1}=20$; $n_{m+1}=21$; $u_{m+1}=8$.\\
Vậy mốt của mẫu số liệu ghép nhóm là
$$M_0=6+\dfrac{37-20}{(37-20)+(37-21)}\cdot(8-6)\approx 7{,}03.$$ 
	\itemch Sai. Vì $x_1; x_2; \ldots; x_{100}$ là mẫu số liệu theo thứ tự không giảm.\\
Tứ phân vị thứ hai của dãy số liệu $x_1; x_2; \ldots; x_{100}$ là $\dfrac{1}{2}\left(x_{50}+x_{51}\right)$.\\
Do $x_{50}$ và $x_{51}$ thuộc nhóm $[6; 8)$ nên $Q_2=6+\dfrac{\dfrac{2 \cdot 100}{4}-32}{37} \cdot(8-6) \approx 6{,}97$.
	\end{itemchoice}
}
\end{ex}
\Closesolutionfile{ans}
\indapan{3}{ans/ans-OTC3-DS}

\Opensolutionfile{ans}[ans/ans-OTC3-KQ]
\caukq
\setcounter{ex}{0}
\begin{ex}%[1D5H1-3]%Câu 1
Khi đo cân nặng của học sinh lớp 11D, y tá lập được bảng số liệu ghép nhóm sau đây:
\begin{center}
\begin{tabular}{|l|c|c|c|c|c|c|}
\hline Cân nặng (kg) & {$[40{,}5; 45{,}5)$} & {$[45{,}5; 50{,}5)$} & {$[50{,}5; 55{,}5)$} & {$[55{,}5; 60{,}5)$} & {$[60{,}5; 65{,}5)$} & {$[65{,}5; 70{,}5)$} \\
\hline Số học sinh & $10$ & $7$ & $16$ & $4$ & $2$ & $3$ \\
\hline
\end{tabular}
\end{center}
Cân nặng trung bình của học sinh lớp 11D là bao nhiêu? (Kết quả làm tròn đến hàng phẩn chục)
\shortans{$51{,}8$}
	\loigiai{
Ta có bảng thống kê
\begin{center}
\begin{tabular}{|l|c|c|c|c|c|c|}
\hline Cân nặng $(\mathrm{kg})$ & {$[40{,}5; 45{,}5)$} & {$[45{,}5; 50{,}5)$} & {$[50{,}5 ; 55{,}5)$} & {$[55{,}5; 60{,}5)$} & {$[60{,}5; 65{,}5)$} & {$[65{,}5 ; 70{,}5)$} \\
\hline Giá trị đại diện & $43$ & $48$ & $53$ & $58$ & $63$ & $68$ \\
\hline Số học sinh & $10$ & $7$ & $16$ & $4$ & $2$ & $3$ \\
\hline
\end{tabular}
\end{center}
Cân nặng trung bình của học sinh là $$\overline{x}=\dfrac{43\cdot 10+48\cdot 7+53\cdot 16+58\cdot 4+63\cdot 2+68\cdot 3}{42}=51{,}8.$$
}
\end{ex}

\begin{ex}%[1D5V1-3]%Câu 2
Thời gian truy cập internet mỗi buổi trưa của một số học sinh được cho trong bảng sau:
\begin{center}
\begin{tabular}{|l|c|c|c|c|c|}
\hline Thời gian (phút) & {$[9{,}5; 12{,}5)$} & {$[12{,}5; 15{,}5)$} & {$[15{,}5 ; 18{,}5)$} & {$[18{,}5; 21{,}5)$} & {$[21{,}5; 24{,}5)$} \\
\hline Số học sinh & $3$ & $12$ & $15$ & $24$ & $2$ \\
\hline
\end{tabular}
\end{center}
Số trung vị của mẫu số liệu trên là bao nhiêu? (Kết quả làm tròn đến hàng phần chục)
\shortans{$18{,}1$}
	\loigiai{
Cỡ mẫu là $n=3+12+15+24+2=56$.\\
Gọi $x_1; \ldots; x_{56}$ là thời gian truy cập internet của $56$ học sinh và giả sử dãy này đã được sắp xếp theo thứ tự tăng dần.\\
Khi đó, số trung vị là $\dfrac{x_{28}+x_{29}}{2}$.\\
Do hai giá trị $x_{28}; x_{29}$ thuộc nhóm $[15{,}5; 18{,}5)$ nên nhóm này chứa trung vị.\\
Do đó $p=3$; $a_3=15{,}5$; $m_3=15$; $m_1+m_2=15$; $a_4-a_3=3$.
$$M_e=15{,}5+\dfrac{\dfrac{56}{2}-15}{15}\cdot 3=18{,}1.$$
}
\end{ex}

\begin{ex}%[1D5V2-3]%Câu 3
Khảo sát thời gian tập thể dục của một số học sinh khối 11 thu được mẫu số liệu ghép nhóm sau:
\begin{center}
\begin{tabular}{|l|c|c|c|c|c|}
\hline Thời gian (phút) & {$[0; 20)$} & {$[20; 40)$} & {$[40; 60)$} & {$[60; 80)$} & {$[80; 100)$} \\
\hline Số học sinh & $5$ & $9$ & $12$ & $10$ & $6$ \\
\hline
\end{tabular}
\end{center}
Tứ phân vị thứ ba của mẫu số liệu là bao nhiêu?
\shortans{$71$}
	\loigiai{
Cỡ mẫu là $n=5+9+12+10+6=42$.\\
Tứ phân vị thứ ba $Q_3$ là $x_{33}$, do nhóm chứa $x_{33}$ là $[60; 80)$, do đó $p=4$; $a_4=60$; $m_1+m_2+m_3=5+9+12=26$; $a_5-a_4=20$.\\
Ta có $Q_3=a_4+\dfrac{\dfrac{3 n}{4}-\left(m_1+m_2+m_3\right)}{m_4}\left(a_5-a_4\right)=60+\dfrac{\dfrac{3.42}{4}-26}{10} \cdot 20=71$.
}
\end{ex}

\begin{ex}%[1D5H1-3]%Câu 4
Khảo sát chiều cao của học sinh lớp $11 \mathrm{C} 1$, thu được bảng số liệu ghép nhóm như sau
\begin{center}
\begin{tabular}{|l|c|c|c|c|c|}
\hline Khoàng chiều cao $(\mathrm{cm})$ & {$[150; 155)$} & {$[155; 160)$} & {$[160; 165)$} & {$[165; 170)$} & {$[170; 175)$} \\
\hline Số học sinh & $7$ & $15$ & $12$ & $8$ & $3$ \\
\hline
\end{tabular}
\end{center}
Chiều cao trung bình của học sinh $11 \mathrm{C} 1$ có dạng $abc{,}d$. Tính $S=a+b+c+d$.
\shortans{$15$}
	\loigiai{
\begin{center}
\begin{tabular}{|l|c|c|c|c|c|}
\hline Khoàng chiều cao $(\mathrm{cm})$ & {$[150; 155)$} & {$[155; 160)$} & {$[160; 165)$} & {$[165; 170)$} & {$[170; 175)$} \\
\hline Giá trị đại diện & $152{,}5$ & $157{,}5$ & $162{,}5$ & $167{,}5$ & $172{,}5$ \\
\hline Số học sinh & $7$ & $15$ & $12$ & $8$ & $3$ \\
\hline
\end{tabular}
\end{center}
Chiều cao trung bình của học sinh lớp 11C1 là
$$\overline{x}=\dfrac{1}{45}(7\cdot 152{,}5+15\cdot 157{,}5+12\cdot 162{,}5+4 \cdot 167{,}5+3\cdot 172{,}5)=160{,}8 \mathrm{~cm}.$$
Khi đó $S=1+6+0+8=15$.
}
\end{ex}

\begin{ex}%[1D5H1-3]%Câu 5
Doanh thu bán hàng trong 20 ngày được lựa chọn ngẫu nhiên của một cửa hàng được ghi lại ở bảng sau (đơn vị: triệu đồng):
\begin{center}
\begin{tabular}{|l|c|c|c|c|c|}
\hline Doanh thu & {$[5; 7)$} & {$[7; 9)$} & {$[9; 11)$} & {$[11; 13)$} & {$[13; 15)$} \\
\hline Số ngày & $2$ & $7$ & $7$ & $3$ & $1$ \\
\hline
\end{tabular}
\end{center}
Doanh số trung bình của cửa hàng trên trong $20$ ngày là bao nhiêu? (Kết quả làm tròn đến hàng phần chục)
\shortans{$9{,}4$}
	\loigiai{
Bảng tần số ghép nhóm theo giá trị đại diện là
\begin{center}
\begin{tabular}{|l|c|c|c|c|c|}
\hline Doanh thu & {$[5; 7)$} & {$[7; 9)$} & {$[9; 11)$} & {$[11; 13)$} & {$[13; 15)$} \\
\hline Giá trị đại diện & $6$ & $8$ & $10$ & $12$ & $14$ \\
\hline Số ngày & $2$ & $7$ & $7$ & $3$ & $1$ \\
\hline
\end{tabular}
\end{center}
Số trung bình $\overline{x}=\dfrac{2\cdot 6+7\cdot 8+7\cdot 10+3\cdot 12+1\cdot 14}{20}=9{,}4$.
}
\end{ex}

\begin{ex}%[1D5V1-4]%Câu 6
Một công ty bất động sản Đất Vàng thực hiện cuộc khảo sát khách hàng xem họ có nhu cầu mua nhà ở mức giá nào để tiến hành dự án xây nhà ở Thăng Long group sắp tới. Kết quả khảo sát $500$ khách hàng được ghi lại ở bảng sau:
\begin{center}
\begin{tabular}{|l|c|c|c|c|c|}
\hline Mức giá (triệu đồng) & {$[10; 14)$} & {$[14; 18)$} & {$[18; 22)$} & {$[22; 26)$} & {$[26; 30)$} \\
\hline Số khách hàng & $75$ & $105$ & $179$ & $96$ & $45$ \\
\hline
\end{tabular}
\end{center}
Công ty bất động sản Đất Vàng nên xây nhà ở mức giá nào để nhiều người có nhu cầu xây nhà?
\shortans{$19{,}9$}
	\loigiai{
Nhóm chứa mốt của mẫu số liệu trên là $[18; 22)$.\\
Do đó $u_m=18$; $n_{m-1}=105$; $n_{m+1}=96$; $u_{m+1}-u_m=22-18=4$.\\
Mốt của mẫu số liệu ghép nhóm là $M_0=18+\dfrac{179-105}{(179-105)+(179-96)} \cdot 4 \approx 19{,}9$.\\
Dựa vào kết quả trên ta có thể dự đoán rằng nếu công ty bất động sản Đất Vàng xây nhà ở mức giá $19{,}9$ triệu đồng$/\mathrm{m}^2$ thì sẽ có nhiều người mua nhất.
}
\end{ex}
\Closesolutionfile{ans}
\indapan{6}{ans/ans-OTC3-KQ}